\documentclass[10pt,a4paper]{article}%
\usepackage[T1]{fontenc}%
\usepackage[utf8]{inputenc}%
\usepackage{lmodern}%
\usepackage{textcomp}%
\usepackage{lastpage}%
\usepackage{times}%
\usepackage{natbib}%
\usepackage{hyperref}%
\usepackage{booktabs}%
\usepackage{amsmath}%
\usepackage{acl}%

\title{Lanka Data: A Four-Field Grammar for Querying Public Data about Sri Lanka}
\author{Nuwan I. Senaratna\\\texttt{nuwans@stanford.edu}}
\date{\today}

\begin{document}
\maketitle

\begin{abstract}
Lanka Data is a software library that provides a single,
uniform interface to public data about Sri Lanka. Rather
than exposing a growing collection of endpoints, methods,
and parameters, it reduces every query to one string of
four positional fields---\emph{what} is measured,
\emph{when}, \emph{where}, and \emph{how} it is
presented---delimited by slashes. The same string serves
unchanged as a Python argument, a command-line argument, a
URL path, and a file path. This paper describes the
motivation for the design, specifies the four-field command
grammar and its single intentional coupling, argues that
the grammar spans the target query space by composition
rather than by accretion, and catalogues the datasets---
census, election, administrative geography, and
hydrology---and their sources that the library exposes.
\end{abstract}

\section{Introduction}

The goal of Lanka Data is ``one API to rule them all'': a
single interface that can express \emph{any} query, rather
than a proliferation of endpoints, methods, libraries, and
parameter sets that each answer one narrow question. Most
data libraries grow by accretion; every new question adds
another function, endpoint, or flag, and no single mental
model survives contact with the result. We wanted the
opposite: a fixed, minimal grammar that a user learns
\emph{once} and can then aim at anything, and that a
non-technical user can read and write without learning to
program. The current domain is public Sri Lankan data---
census measurements, election results, and administrative
geography---but nothing in the grammar is specific to Sri
Lanka: \emph{what}, \emph{when}, \emph{where}, and
\emph{how} are the dimensions of essentially any factual
query about the world.

\section{The Command Grammar}

The public interface is a single string parsed into four
positional fields, delimited by slashes:

\begin{center}
\texttt{What / When / Where / How}
\end{center}

There is no secondary configuration surface: no options
object, no builder API, no config files. Each field is an
independent axis of the query, and the value chosen for one
field does not constrain the valid values of another, with a
single intentional coupling described below.

\subsection{What --- the measurement}

\textbf{What} identifies the quantity being retrieved. It
is a measurement, independent of time, region, and
presentation. Measurements include census quantities and
election results. A reserved keyword denotes the
\emph{absence} of a measurement: a request for region
geometry with no data bound to it, used for base maps and
for inspecting boundary changes in isolation. Because
\textbf{What} encodes only the measurement type, a single
measurement is reused across every combination of the
other three fields without expanding the vocabulary.

\subsubsection*{census-population (29 values)}
AgeGroup, AgriOccupations, Attendance, Dependency,
Disability, Economy, Education, Employment, Enrollment,
Ethnicity, Fertility, Gender, Growth, Inactive, Industry,
Laborforce, Literacy, Marital, Migration,
NonAgriEmployment, NotAttending, Occupations, Population,
RelationshipToHead, Religion, Sectoral, Sectors,
Speaking, Unemployment

\subsubsection*{census-housing (22 values)}
Communication, ConstructionYear, Electricity, Floor, Fuel,
Housing, Informal, Lighting, Materials, Occupancy,
Ownership, Persons, Quarters, Roof, Rooms, Structure,
Tenure, Toilet, Unit, Walls, Waste, Water

\subsubsection*{election (6 values)}
Local, LocalSummary, Parliamentary, ParliamentarySummary,
Presidential, PresidentialSummary

\subsubsection*{rivers (2 values)}
Catchment, RiverLen

\subsection{When --- the observation time}

\textbf{When} binds the measurement to the point or
interval in time at which it happened. It accepts either a
single year or an interval. An interval is a single value,
not two concatenated queries.

Supported formats:
\begin{itemize}
\item Single year: e.g., \texttt{2024}
\item Interval: e.g., \texttt{2012-2024}
\item If exact year unavailable, closest available data
is returned
\end{itemize}

\subsection{Where --- the region}

\textbf{Where} identifies the region under measurement: its
identity within the administrative hierarchy, not merely a
coordinate. It carries the most syntax of the four fields
because it is the axis along which real queries vary most.
The supported forms are a single region, resolution into
child regions of a given type, an explicit set of named
regions, a contiguous range between two endpoints, and a
historical boundary variant that selects a region's
boundaries as they existed before a given boundary
redesign. Observation time and boundary epoch are kept as
separate values so that counts are not misattributed to the
wrong geometry.

Supported syntax patterns:
\begin{itemize}
\item \texttt{<region\_id>}: Single region (e.g., LK)
\item \texttt{<region\_id>:<type>}: Child regions
(e.g., LK:district)
\item \texttt{<r1>,<r2>}: Multiple regions
(e.g., LK-1,LK-2)
\item \texttt{<r1>...<r2>}: Range between endpoints
(e.g., LK-1...LK-2)
\item \texttt{<region>@<distance>}: Regions within
distance of same type
\end{itemize}

\subsection{How --- the presentation}

\textbf{How} specifies the output representation. The same
measurement can be emitted as a map, a chart, or raw data
without any change to \textbf{What}. Because presentation
is separated from measurement, \textbf{How} carries its own
modifier grammar: rankings of a category, shares,
differences between observations, and diversity indices are
transformations applied to the same underlying values. This
is the single intentional coupling of the grammar:
change-based modifiers are valid only when \textbf{When}
supplies an interval.

Supported visualization bases (24 total):
BarChart, BivariateMap, BubbleMap, BumpChart, CSV,
ChartSpec, GeoJSON, HexMap, Histogram, JSON, LineChart,
Map, Parquet, PieChart, QuadrantChart, ScatterPlot,
SquareMap, StackedBarChart, TSV, TreeMap, TriangleMap,
UnitHexMap, UnitSquareMap, UnitTriangleMap

Supported modifiers (12 total):
1st, Top, 2nd, 3rd, Bottom, 1stPct, 2ndPct, Change,
Top3, Diversity, DiversityPew

\section{Datasets and Sources}

The \textbf{What} vocabulary is organised into four data
categories. \emph{census-population} exposes measurements
such as ethnicity, religion, age group, education,
employment, and literacy; \emph{census-housing} exposes
housing characteristics such as construction materials,
water, lighting, and toilet facilities;
\emph{election} exposes presidential, parliamentary, and
local government results and summaries; and \emph{rivers}
exposes river length and catchment statistics. Each
response carries the provenance of the sources below.

\subsection{Census of Population and Housing}

The census-population category exposes measurements such as
ethnicity, religion, age group, education, employment, and
literacy. The census-housing category exposes housing
characteristics such as construction materials, water,
lighting, and toilet facilities. Data is available from the
2001, 2012, and 2024 census cycles \citep{census_2024}.
All census data is sourced from the Department of Census
and Statistics.

\subsection{Elections}

The election category exposes presidential, parliamentary,
and local government results and summaries across multiple
election years. This includes vote tallies, seat
distributions, and electoral statistics at national,
district, and constituency levels. All election data is
sourced from the \citet{elections_commission}.

\subsection{Administrative Geography}

The administrative boundary data includes both current and
historical boundary epochs, allowing queries to be anchored
to the boundaries as they existed at the time of
observation. This includes provinces, districts, divisional
secretariat divisions, and Grama Niladhari divisions across
different time periods. Data is sourced from the
\citet{dcs_sri_lanka}.

\subsection{Rivers}

The rivers category exposes river length and catchment
statistics for major waterways in Sri Lanka \citep{hydrorivers}. This data includes geometric
and hydrological properties of river systems. Data is
sourced from HydroRIVERS via the \texttt{lk\_rivers}
project.


\bibliographystyle{acl_natbib}
\bibliography{lanka_data}
\end{document}